\documentclass[11pt,a4paper]{report}

% Packages
\usepackage[margin=0.9in]{geometry}
\usepackage{fancyhdr}
\usepackage{titlesec}
\usepackage{booktabs}
\usepackage{longtable}
\usepackage{array}
\usepackage{amsmath}
\usepackage{enumitem}
\usepackage{hyperref}
\usepackage{xcolor}
\usepackage{graphicx}
\usepackage{float}
\usepackage{verbatim}
\usepackage{listings}
\usepackage{tikz}
\usetikzlibrary{shapes.geometric, arrows.meta, positioning, calc, fit, backgrounds}

% Colors
\definecolor{urteg}{RGB}{0,128,128}
\definecolor{arcadiablue}{RGB}{0,82,147}
\definecolor{capellagreen}{RGB}{0,128,0}
\definecolor{sectblue}{RGB}{0,51,102}

% Listings for SysML/Arcadia textual schemas
\lstdefinestyle{sysmlstyle}{
    basicstyle=\ttfamily\footnotesize,
    breaklines=true,
    frame=single,
    backgroundcolor=\color{gray!8},
    keywordstyle=\color{arcadiablue}\bfseries,
    commentstyle=\color{gray},
    numbers=left,
    numberstyle=\tiny\color{gray}
}

% Header/Footer
\pagestyle{fancy}
\fancyhf{}
\fancyhead[L]{\small URTE Arcadia/Capella MBSE Extension | Confidential}
\fancyhead[R]{\small June 2026 | v1.0}
\fancyfoot[C]{\thepage}
\renewcommand{\headrulewidth}{0.4pt}

% Section formatting
\titleformat{\chapter}{\Large\bfseries\color{sectblue}}{\thechapter.}{0.5em}{}
\titleformat{\section}{\large\bfseries\color{arcadiablue}}{\thesection}{0.5em}{}
\titleformat{\subsection}{\normalsize\bfseries\color{urteg}}{\thesubsection}{0.5em}{}

% Title
\title{\textbf{URTE Arcadia / Capella MBSE Extension}\\[0.3cm]
\large SysML Viewpoints and Extended Schemas\\[0.2cm]
\normalsize Compliant with \textbf{Arcadia} Method and \textbf{Capella} Tooling\\[0.4cm]
for the\\
\textbf{Universal Regeneration Therapy Ecosystem (URTE)}}
\author{
    \textbf{Prepared by:}\\[0.1cm]
    Grok xAI Research \& UAB Saptera\\[0.1cm]
    Kristupas Gelzinis \& Olek Suchodolski\\[0.3cm]
    \textit{Extension of URTE MBSE Document (June 2026)}\\[0.1cm]
    \textit{Based on URTE Architecture Document, Investment Prospectus,}\\
    \textit{Infrastructure Model, and Situation Room Briefing}
}
\date{June 2026}

\begin{document}

\maketitle

\begin{abstract}
\noindent This document extends the URTE MBSE baseline with detailed \textbf{Arcadia}-compliant viewpoints and \textbf{SysML}-like schemas suitable for direct implementation in the \textbf{Capella} open-source MBSE tool.

It applies the full Arcadia method (Operational Analysis, System Analysis, Logical Architecture, Physical Architecture, and EPBS) to the Universal Regeneration Therapy Ecosystem. All schemas are provided in textual/SysML-inspired notation that can be recreated in Capella using standard Arcadia viewpoints, allowing traceability from operational needs through to physical deployment while preserving the scale-invariant, TRL Pull, and ethical-by-design principles of URTE.
\end{abstract}

\tableofcontents
\newpage

%==============================================================================
\chapter{Introduction to Arcadia Method for URTE}
%==============================================================================

\section{Why Arcadia / Capella for URTE?}

The URTE system is a complex, multi-scale, socio-technical System-of-Systems involving:
- Biological and synthetic components (bionanorobots, cell/gene therapies)
- AI-intensive software (digital twins, generative design, RL optimizers)
- Physical deployment across clinical, terrestrial, and space environments
- Strong ethical, regulatory, and planetary-boundary governance constraints

\textbf{Arcadia} (developed by Thales) provides a structured, viewpoint-based methodology that is particularly well-suited because it emphasizes:
- Clear separation between \textbf{what} the system must do (Operational & System Analysis) and \textbf{how} it is realized (Logical & Physical Architecture)
- Strong support for \textbf{functional chains} and \textbf{data flows} across scales — critical for URTE’s cross-scale coupling
- Built-in traceability and impact analysis — essential for TRL Pull mechanism and regulatory compliance
- Native support in the free \textbf{Capella} tool, enabling collaborative MBSE without heavy licensing costs

This extension document maps the existing URTE 6-layer architecture and biomathematical models into standard Arcadia viewpoints, providing ready-to-use schemas for Capella implementation.

\section{Arcadia Viewpoint Structure Applied to URTE}

\begin{tabular}{>{\raggedright\arraybackslash}p{3.5cm} p{10cm}}
\toprule
\textbf{Arcadia Viewpoint} & \textbf{URTE Focus \& Key Artifacts} \\
\midrule
\textbf{Operational Analysis (OA)} & Operational Capabilities (e.g., ``Regenerate damaged tissue at subcellular scale'', ``Restore ecosystem resilience at landscape scale''), Operational Entities (Clinicians, Ecologists, Space Agencies, Ethics Boards), Operational Activities and Scenarios \\
\textbf{System Analysis (SA)} & System Capabilities, System Functions (allocated from the 6 layers), Functional Chains (Intervention Lifecycle, TRL Pull Flywheel, Cross-Scale Learning), System Actors, Data Model (multi-scale state vectors $u_s$) \\
\textbf{Logical Architecture (LA)} & Logical Components (the 6 URTE layers as reusable Logical Components with Ports \& Interfaces), Logical Functions, Logical Data Flows, Allocation of Functions to Components \\
\textbf{Physical Architecture (PA)} & Physical Nodes (Ground Mission Control, Clinical Deployment Hubs, Field Kits, MedInt Satellite Constellation, Orbital/Lunar Precursors), Physical Components, Deployment of Logical Components, Physical Links \\
\textbf{EPBS} & Configuration Items, Versioned Builds, Safety/Security Classification per Component \\
\bottomrule
\end{tabular}

%==============================================================================
\chapter{Operational Analysis (OA) Viewpoint}
%==============================================================================

\section{Operational Capabilities}

\begin{enumerate}[label=OC-\arabic*]
    \item Regenerate living systems at molecular/cellular scale with precision and safety.
    \item Restore ecosystem functions and resilience at landscape/regional scale.
    \item Enable controlled planetary regeneration precursors and bio-regen + ISRU for space habitats.
    \item Accelerate technology maturation across domains via TRL Pull mechanism.
    \item Enforce ethical, planetary-boundary, and governance constraints on all interventions.
\end{enumerate}

\section{Operational Entities \& Actors}

\begin{itemize}
    \item \textbf{Clinician / Patient Entity}
    \item \textbf{Ecologist / Restoration Practitioner}
    \item \textbf{Space Agency / ISRU Engineer}
    \item \textbf{Multi-Stakeholder Ethics Board}
    \item \textbf{Regulatory Authority (FDA/EMA/UN)}
    \item \textbf{General Public / Civil Society}
\end{itemize}

\section{Operational Architecture Diagram (Textual SysML-inspired)}

\begin{lstlisting}[style=sysmlstyle, caption={Operational Architecture (OA) -- Textual Representation for Capella}]
Operational Architecture [OA]
├── Operational Entities
│   ├── Clinician/Patient
│   ├── Ecologist/Restoration Team
│   ├── Space Agency/ISRU Team
│   ├── Ethics Board (Multi-Stakeholder)
│   └── Regulatory Authority
├── Operational Activities
│   ├── OA.1 Perform Subcellular Repair
│   ├── OA.2 Execute Landscape Restoration
│   ├── OA.3 Conduct Mars/Moon Analog Campaign
│   ├── OA.4 Authorize High-Impact Intervention (Ethics Gate)
│   └── OA.5 Monitor Planetary Boundary Indicators
├── Operational Capabilities
│   ├── OC.1 Precision Multi-Scale Regeneration
│   ├── OC.2 TRL Pull Acceleration
│   └── OC.3 Ethical & Planetary Guardrail Enforcement
└── Operational Scenarios
    ├── Scenario 1: First-in-Human Bionanorobot Therapy (with sandbox progression)
    └── Scenario 2: Cross-Domain Learning -- Clinical success improves ecosystem model
\end{lstlisting}

%==============================================================================
\chapter{System Analysis (SA) Viewpoint}
%==============================================================================

\section{System Capabilities}

\begin{enumerate}[label=SC-\arabic*]
    \item SC.1 Provide unified multi-scale digital twin simulation (molecular $\leftrightarrow$ planetary)
    \item SC.2 Generate and optimize intervention protocols under guardrails
    \item SC.3 Execute progressive sandboxed physical interventions
    \item SC.4 Enable closed-loop cross-domain learning and TRL Pull
    \item SC.5 Maintain full traceability, auditability, and ethical oversight
\end{enumerate}

\section{System Functions (Derived from URTE 6 Layers)}

\begin{longtable}{>{\raggedright\arraybackslash}p{2.2cm} p{8cm} p{3.5cm}}
\toprule
\textbf{Function ID} & \textbf{System Function Description} & \textbf{Allocated to Layer} \\
\midrule
\endfirsthead
\toprule
\textbf{Function ID} & \textbf{System Function Description} & \textbf{Allocated to Layer} \\
\midrule
\endhead
\bottomrule
\endfoot
F.1.1 & Authorize intervention via multi-stakeholder ethics board & Layer 1 \\
F.1.2 & Enforce planetary-boundary mathematical guardrails & Layer 1 + 2 \\
F.2.1 & Run Physics-Informed Neural Network (PINN) multi-scale simulation & Layer 2 \\
F.2.2 & Generate candidate intervention protocols (generative design) & Layer 2 \\
F.2.3 & Optimize interventions via stochastic optimal control (J(u)) & Layer 2 \\
F.2.4 & Perform uncertainty quantification & tipping-point early warning & Layer 2 \\
F.3.1 & Execute high-TRL clinical regen modules (cell/gene, bioprinting) & Layer 3 \\
F.4.1 & Execute eco-planetary regen protocols & Layer 4 \\
F.5.1 & Deploy bionanorobot swarms with kill-switches & containment & Layer 5 \\
F.5.2 & Collect real-time telemetry from all scales & Layer 5 \\
F.6.1 & Update digital twin models from outcome data & Layer 6 \\
F.6.2 & Perform cross-domain learning (clinical $\leftrightarrow$ ecological) & Layer 6 \\
\bottomrule
\end{longtable}

\section{Key Functional Chain: Intervention Lifecycle with TRL Pull (Activity Diagram - Textual)}

\begin{lstlisting}[style=sysmlstyle, caption={Functional Chain: End-to-End Intervention Lifecycle + TRL Pull}]
Functional Chain: Intervention Lifecycle with TRL Pull (SA)
Start
  --> [F.1.1] Ethics Board Authorization Gate
  --> [F.2.2] Generative Protocol Designer (Layer 2)
  --> [F.2.3] RL Optimizer with Guardrails (R, planetary boundaries)
  --> Decision: High-Impact? --> [F.1.1] Ethics Veto possible
  --> [F.5.1] Progressive Sandbox Deployment (in-vitro -> animal -> controlled pilot)
  --> [F.5.2] Real-time Telemetry Collection
  --> [F.6.1] Digital Twin Model Update
  --> [F.6.2] Cross-Domain Learning --> TRL Pull to lower-maturity modules
  --> End (Outcome logged + Audit trail complete)
\end{lstlisting}

\section{Data Model -- Multi-Scale State Vector (SysML Block)}

\begin{lstlisting}[style=sysmlstyle, caption={Data Model: Multi-Scale State Vector (SysML Block Definition)}]
Block: MultiScaleStateVector
  Attributes:
    + scaleLevel: ScaleEnum {Molecular, Cellular, Tissue, Organ, Ecosystem, Regional, Planetary}
    + u_s: StateVector          // concentrations, stress, biomass, CO2, etc.
    + lambda_dom: Real          // dominant eigenvalue for resilience R
    + variance: Real            // early-warning indicator
    + lag1_autocorr: Real
    + skewness: Real
    + kurtosis: Real
  Operations:
    + computeResilience(): Real     // R = -lambda_dom
    + detectTippingPoint(): Boolean
  Constraints:
    {planetaryBoundaryGuardrails enforced on all u_s updates}
\end{lstlisting}

%==============================================================================
\chapter{Logical Architecture (LA) Viewpoint}
%==============================================================================

\section{Logical Components (Mapping URTE 6 Layers)}

The six URTE layers are modeled as reusable **Logical Components** with well-defined ports and interfaces. This enables strong modularity and traceability.

\begin{lstlisting}[style=sysmlstyle, caption={Logical Architecture (LA) -- URTE Layers as Logical Components}]
Logical Architecture [LA]
├── Logical Component: CommandAndGovernance (Layer 1)
│   Ports:
│     - EthicsVetoPort (in/out)
│     - RegulatoryInterfacePort
│     - AuditLogPort
│   Allocated Functions: F.1.1, F.1.2
│
├── Logical Component: AIOrchestrationAndMultiScaleSimulation (Layer 2)
│   Ports:
│     - TwinInputPort (sensor data, u_s)
│     - ProtocolOutputPort (candidate interventions)
│     - GuardrailEnforcementPort
│   Allocated Functions: F.2.1 -- F.2.4
│
├── Logical Component: BiologicalRegenSubsystem (Layer 3)
│   Ports:
│     - ClinicalDataPort
│     - TherapyExecutionPort
│   Allocated Functions: F.3.1
│   Note: Primary TRL Pull revenue & data source (TRL 6-9)
│
├── Logical Component: EcoPlanetaryRegenSubsystem (Layer 4)
│   Ports:
│     - EcoTelemetryPort
│     - RestorationProtocolPort
│   Allocated Functions: F.4.1
│
├── Logical Component: SensingActuationAndDeployment (Layer 5)
│   Ports:
│     - TelemetryOutPort
│     - InterventionCommandPort
│     - KillSwitchPort (safety)
│   Allocated Functions: F.5.1, F.5.2
│   Sub-Components:
│     - BionanorobotSwarm (with velocity regime models)
│     - IoTSensorNetwork
│     - BioreactorBioprinterFleet
│
└── Logical Component: ContinuousFeedbackAndLearningLoop (Layer 6)
    Ports:
      - ModelUpdatePort
      - CrossDomainLearningPort
    Allocated Functions: F.6.1, F.6.2
    Note: Implements Virtuous Evidence Flywheel + TRL Pull
\end{lstlisting}

\section{Logical Internal Block Diagram (IBD) -- Simplified Textual}

\begin{lstlisting}[style=sysmlstyle, caption={Logical IBD: Cross-Layer Data Flow (Arcadia Style)}]
Logical IBD: URTE Core Data Flow
[CommandAndGovernance] --EthicsVeto--> [AIOrchestration...]
[AIOrchestration...] --Protocol--> [SensingActuation...]
[SensingActuation...] --Telemetry--> [ContinuousFeedback...]
[ContinuousFeedback...] --ModelUpdate--> [AIOrchestration...]
[BiologicalRegen] --ClinicalOutcomeData--> [ContinuousFeedback...]  // TRL Pull source
[EcoPlanetaryRegen] --EcoOutcomeData--> [ContinuousFeedback...]
\end{lstlisting}

%==============================================================================
\chapter{Physical Architecture (PA) Viewpoint}
%==============================================================================

\section{Physical Nodes \& Deployment}

\begin{lstlisting}[style=sysmlstyle, caption={Physical Architecture (PA) -- Nodes and Deployment}]
Physical Architecture [PA]
├── Physical Node: GroundMissionControlCenter
│   Hosts: CommandAndGovernance (Layer 1)
│          AIOrchestrationAndMultiScaleSimulation (Layer 2)
│          ContinuousFeedbackAndLearningLoop (Layer 6)
│   Links: High-bandwidth to Clinical Hubs and Satellite Ground Stations
│
├── Physical Node: ClinicalDeploymentHub
│   Hosts: BiologicalRegenSubsystem (Layer 3) + parts of Layer 5
│   Components: Bioprinters, Robotic Surgical Systems, Local Telemetry Cache
│
├── Physical Node: FieldEcoRestorationKit (mobile/edge)
│   Hosts: parts of EcoPlanetaryRegenSubsystem + SensingActuation
│   Components: Drone Swarms, IoT Gateways, Portable Bioreactors
│
├── Physical Node: MedIntSatelliteConstellation (Ecological Measurement)
│   Hosts: parts of Layer 5 (remote sensing + telemetry relay)
│   Links: to GroundMissionControl and SpaceStationGovernanceNode
│
└── Physical Node: SpaceStationGovernanceNode (Phase 3+)
    Hosts: Governance oversight for MedInt constellation + precursor bio-regen payloads
    Links: to Orbital/Lunar Bio-Regen + ISRU Modules
\end{lstlisting}

\section{Physical Deployment Diagram (Textual)}

\begin{lstlisting}[style=sysmlstyle, caption={Physical Deployment: URTE Components to Nodes}]
Deployment:
  CommandAndGovernance          --> GroundMissionControlCenter
  AIOrchestration...            --> GroundMissionControlCenter
  BiologicalRegenSubsystem      --> ClinicalDeploymentHub
  EcoPlanetaryRegenSubsystem    --> FieldEcoRestorationKit + GroundMissionControlCenter
  SensingActuation...           --> ClinicalDeploymentHub + FieldEcoRestorationKit + MedIntSatellites
  ContinuousFeedback...         --> GroundMissionControlCenter
\end{lstlisting}

%==============================================================================
\chapter{EPBS and Configuration Management}
%==============================================================================

\textbf{End-Product Breakdown Structure (EPBS)} for URTE Phase 1 (recommended Capella structure):

\begin{itemize}
    \item \textbf{URTE Core Platform v1.0} (Configuration Item)
          \begin{itemize}
              \item Digital Twin Engine (PINN-based)
              \item AI Orchestration Module
              \item Governance \& Ethics Dashboard
              \item Bionanorobot Sandbox Kit (v0.8)
          \end{itemize}
    \item \textbf{Clinical Pilot Deployment Package}
    \item \textbf{Terrestrial Eco-Restoration Pilot Package}
    \item \textbf{Data Standards \& API Library} (open source part)
\end{itemize}

Each Configuration Item carries safety classification, TRL level, and traceability to requirements.

%==============================================================================
\chapter{Cross-Viewpoint Traceability \& Capella Implementation Guidance}
%==============================================================================

\section{Traceability Matrix (Excerpt)}

\begin{tabular}{lll}
\toprule
\textbf{Arcadia Element} & \textbf{URTE MBSE Requirement} & \textbf{TRL / Phase} \\
\midrule
Operational Capability OC.1 & FR-01, SR-01 & Phase 1--3 \\
Functional Chain "Intervention Lifecycle + TRL Pull" & FR-04, NFR-03 & All Phases \\
Logical Component "AIOrchestration..." & FR-01, FR-02 & Phase 1 \\
Physical Node "MedIntSatelliteConstellation" & Phase 3 roadmap & Phase 3 \\
\bottomrule
\end{tabular}

\section{Recommended Capella Project Structure}

\begin{enumerate}
    \item Create new Capella project: \texttt{URTE\_MBSE\_v1}
    \item Use standard Arcadia viewpoints in order:
          \begin{enumerate}
              \item Operational Analysis
              \item System Analysis
              \item Logical Architecture
              \item Physical Architecture
              \item EPBS
          \end{enumerate}
    \item Define reusable libraries for:
          \begin{itemize}
              \item MultiScaleStateVector (Data)
              \item BionanorobotSwarm (Logical + Physical)
              \item GuardrailConstraints (Constraints)
          \end{itemize}
    \item Use Capella’s built-in traceability matrix and impact analysis views for TRL Pull and regulatory evidence.
    \item Export diagrams as images or SVG for inclusion in documentation (or use Capella’s HTML/Word export).
\end{enumerate}

%==============================================================================
\chapter{Conclusion and Next Steps}
%==============================================================================

This Arcadia/Capella extension provides a complete set of viewpoints and textual SysML schemas that faithfully represent the URTE system while being directly implementable in the free Capella tool.

\textbf{Recommended immediate next steps:}
\begin{enumerate}
    \item Import the Logical Components and Functional Chains into Capella.
    \item Create the first version of the multi-scale data model as a shared library.
    \item Define initial Operational Scenarios for the four Phase 1 pilots.
    \item Run Capella’s validation rules to ensure consistency across viewpoints.
    \item Generate documentation exports for integration into the main URTE MBSE Document.
\end{enumerate}

The combination of the original URTE MBSE Document + this Arcadia extension gives a full, tool-ready MBSE baseline aligned with both academic best practices (SEBoK/SWEBOK) and industrial methodology (Arcadia/Capella).

\vspace{0.8cm}
\begin{center}
\textit{--- End of URTE Arcadia / Capella MBSE Extension v1.0 ---}
\end{center}

%==============================================================================
\appendix
%==============================================================================

\chapter{Quick Reference: Arcadia Viewpoint Mapping to URTE Layers}

\begin{tabular}{p{4cm}p{5cm}p{5cm}}
\toprule
\textbf{Arcadia View} & \textbf{URTE Layer Focus} & \textbf{Key SysML Element} \\
\midrule
Operational Analysis & Stakeholder needs across clinical, eco, space & Operational Capabilities + Scenarios \\
System Analysis & What the system must achieve (6-layer functions) & System Functions + Functional Chains \\
Logical Architecture & How the system is structured (reusable components) & Logical Components (6 layers) + Ports \\
Physical Architecture & Where and how it is deployed (nodes) & Physical Nodes + Deployment \\
EPBS & Configuration \& build management & Configuration Items \\
\bottomrule
\end{tabular}

\chapter{Selected Extended ASCII SysML Diagrams}

\section{Block Definition Diagram (BDD) -- URTE Core Logical Components}

\begin{lstlisting}[style=sysmlstyle]
BDD: URTE_Logical_Components
block CommandAndGovernance
block AIOrchestrationAndMultiScaleSimulation
block BiologicalRegenSubsystem
block EcoPlanetaryRegenSubsystem
block SensingActuationAndDeployment {
    block BionanorobotSwarm
}
block ContinuousFeedbackAndLearningLoop

AIOrchestrationAndMultiScaleSimulation --> BiologicalRegenSubsystem : <<allocate>>
AIOrchestrationAndMultiScaleSimulation --> EcoPlanetaryRegenSubsystem : <<allocate>>
SensingActuationAndDeployment --> BionanorobotSwarm : composition
\end{lstlisting}

\section{State Machine Diagram -- Intervention Authorization \& Sandbox Progression}

\begin{lstlisting}[style=sysmlstyle]
StateMachine: InterventionLifecycle
[*] --> EthicsReview
EthicsReview --> Authorized : ethicsBoard.approve()
EthicsReview --> Rejected : ethicsBoard.reject()
Authorized --> InVitroSandbox
InVitroSandbox --> SmallAnimalSandbox : validation.pass()
InVitroSandbox --> Paused : safetyTrigger()
SmallAnimalSandbox --> ControlledPilot : validation.pass()
ControlledPilot --> OperationalDeployment : regulatory.approval()
OperationalDeployment --> Completed
Paused --> EthicsReview : rootCauseResolved()
\end{lstlisting}

\end{document}
